\documentclass[11pt,a4paper]{article}
% Compile with XeLaTeX (fontspec and xeCJK).
\usepackage{fontspec}
\setmainfont{LibertinusSerif}[Extension=.otf,UprightFont=*-Regular,BoldFont=*-Bold,ItalicFont=*-Italic,BoldItalicFont=*-BoldItalic]
\setsansfont{LibertinusSans}[Extension=.otf,UprightFont=*-Regular,BoldFont=*-Bold,ItalicFont=*-Italic]
\usepackage{xeCJK}
\setCJKmainfont{Hiragino Mincho ProN}
\xeCJKDeclareCharClass{Default}{"00B0, "00B7, "2013, "2014, "2018 -> "201F, "2026}
\usepackage[french]{babel}
\usepackage[a4paper,margin=2.5cm,top=2.3cm]{geometry}
\usepackage{graphicx}
\usepackage{xcolor}
\usepackage{titlesec}
\usepackage{fancyhdr}
\usepackage{enumitem}
\usepackage{xurl}
\usepackage[hidelinks]{hyperref}
\definecolor{accent}{RGB}{31,73,125}
\definecolor{muted}{RGB}{93,104,119}
\hypersetup{colorlinks=true,urlcolor=accent,linkcolor=accent,pdfauthor={Romain Fernex},pdftitle={Rapport sur l’accord de libre-échange UE–Corée du Sud: discussion sur l’histoire et l’avenir de la politique commerciale européenne en Asie}}
\urlstyle{same}
\setlist[itemize]{topsep=.25em,itemsep=.15em,parsep=0pt,leftmargin=1.4em}
\setcounter{secnumdepth}{0}
\titleformat{\section}{\large\bfseries\color{accent}}{}{0em}{}
\titlespacing*{\section}{0pt}{1.6em}{.6em}
\setlength{\parindent}{1.2em}
\setlength{\parskip}{0.35em}
\pagestyle{fancy}
\fancyhf{}
\fancyhead[L]{\small\color{muted}Romain Fernex}
\fancyhead[R]{\small\color{muted}\itshape Accord de libre-échange UE–Corée du Sud}
\fancyfoot[C]{\small\color{muted}\thepage}
\renewcommand{\headrulewidth}{0pt}
\newcommand{\partbreak}{\par\begin{center}\color{muted}$\ast$\quad$\ast$\quad$\ast$\end{center}\par}


\begin{document}
\thispagestyle{empty}
\noindent{\small\sffamily\color{accent}\MakeUppercase{Rapport · Politique commerciale européenne}}\par\medskip
\noindent{\LARGE\bfseries Rapport sur l’accord de libre-échange UE–Corée du Sud\par}
\smallskip\noindent{\large\itshape Discussion sur l’histoire et l’avenir de la politique commerciale européenne en Asie\par}
\medskip
\noindent{\color{accent}\rule{\linewidth}{0.8pt}}\par\smallskip
\noindent{\small\color{muted}Romain Fernex · 2023 · Rapport en français}\par
\smallskip\noindent{\small\color{muted}\emph{Article publié en français sur :} \url{https://www.callonu-escpeurope.com/articles/l-accord-de-libre-échange-ue-corée-du-sud-discussion-sur-l-histoire-et-l-avenir-de-la-politique-commerciale-européenne-en-asie}}\par\bigskip

Le 2 novembre 2022, William Yuen Yee, assistant chercheur pour le «~Columbia-Harvard China and the World program~», publiait, pour le journal The Diplomat, un article portant sur les grandes difficultés rencontrées par l'UE pour obtenir la signature d'un nouvel accord de libre-échange (ALE) avec l'Indonésie. Il y pointe les déboires d'une UE prompte à rappeler sa volonté d'étendre ses relations économiques avec l'Asie\textbf{{[}1{]}} dans un contexte international toujours plus tendu, en particulier avec les disruptions de supply chain qui se multiplient, mais qui se retrouve souvent freinée par sa volonté d'imposer à ses partenaires des normes environnementales qui sont vues comme des contraintes insurmontables par la grande majorité des États de la zone Indopacifique. De manière générale, la présence économique de l'Europe en Asie a été relativement tardive et marquée dès ses débuts par une ambition notable de défendre des normes strictes d'un point de vue social et environnemental. De fait, c'est un des aspects déterminants de la politique bilatérale européenne en termes de commerce international. Pour mieux en comprendre la spécificité et les débats qui en découlent, cet article se propose de revenir sur le tout premier accord de libre-échange de l'UE en Asie, précurseur de la politique commerciale européenne en Asie et plus généralement de sa conception des relations économiques internationales sur le plan bilatéral. S'il est impossible de résumer une décennie de commerce UE-Asie en un article, nous tenterons de voir en quoi l'ALE signé en 2007 entre l'Union Européenne et la Corée du Sud est le reflet de ce qui fait la particularité et l'ambiguïté des relations économiques entre l'UE et l'Asie.

Nous nous intéresserons tout d'abord à la composition de l'accord en lui-même en nous concentrant sur ces chapitres les plus novateurs. Puis nous verrons l'impact économique de celui-ci sur les relations entre la Corée du Sud et l'Europe. Enfin nous nous pencherons sur les différentes controverses qui entourent celui-ci et étendront notre réflexion à certains des ALEs signés par l'UE en Asie après celui-ci.

\partbreak

Comme dans tout ALE, plusieurs clauses prévues dans le chapitre 2 de l'accord visent à la suppression quasi totale des barrières tarifaires existantes -- de l'ordre de 99\% pour les produits industriels - , à l'exception de certains secteurs spécifiques comme la pêche ou l'agriculture. Plus intéressant, ce chapitre contient aussi des provisions concernant les barrières non tarifaires i.e. les obstacles au commerce autres que les droits de douanes \textbf{{[}2{]},} comme l'existence de certaines normes spécifiques à un pays. Par exemple, dans le secteur automobile, l'accord prévoit l'adoption des standards de sécurité internationaux par la Corée du Sud dans le but d'harmoniser les normes entre les deux pays. Un système de «~règlement des différends accéléré~» est également mise en place pour permettre un arbitrage plus rapide en cas de litige. Cette attention aux barrières non tarifaires se retrouve dans le chapitre 4 sur les barrières techniques au commerce ou encore dans le chapitre 6 sur les douanes et la facilitation des échanges. Enfin, le chapitre 10 sur la propriété intellectuelle présente un réel intérêt pour l'Union Européenne, bien plus que pour la Corée, et vise à garantir, entre autres, le respect des indications lié à l'origine géographique de certains produits, notamment issus de l'agro-alimentaire. Ainsi les appellations de style AOP etc.. bénéficient d'un niveau de protection plus élevé. Dans ce même esprit, l'accord comprend un principe de «~rule of origin~» qui impose un seuil minimum de 45\% de composants fabriqués dans le pays d'origine pour que le produit final exporté soit considéré comme en étant originaire.

Au-delà de ces considérations principalement commerciales, un chapitre est particulièrement novateur~: le chapitre 13 intitulé commerce et développement durable. A l'image des accords mis en place par l'UE pour l'intégration des pays d'Europe de l'Est, la signature accord se voit conditionnée à l'acceptation d'un nombre conséquents de normes environnementales et sociales. Des structures institutionnelles sont mises en place, dans le cadre de l'ALE, par chaque partie, avec la présence d'acteurs de la société civil, pour s'assurer que ces engagements sont respectés. Ainsi, pour ce qui est du droit du travail, les 2 parties s'engagent à suivre scrupuleusement les règles éditées par l'organisation internationale du travail (OIT), alors même que la Corée du Sud n'avait ratifié que 4 des 8 grands axes de l'OIT \textbf{{[}3{]}.} Sur le plan environnemental les deux parties promettent, entre autres, de mettre en place concrètement les moyens de tenir leurs engagements vis-à-vis des différents accords environnementaux qu'ils ont signé -- dans le cadre des COPs par exemple .

Cet ensemble de mesures, bien qu'encore limitées, est une première pour l'époque et l'accord de libre échangé négocié par les Etats-Unis avec la Corée du Sud au même moment n'inclut aucune provision de ce type.

\partbreak

Mais qu'en est-il de l'impact économique et social de cet accord pour les deux parties~? Pour l'Union Européenne, dont le volume total des exportations en direction de la Corée augmente de 70\% de 2010 à 2020, le pari se révèle gagnant. Une hausse qui ne porte pas que sur les biens manufacturés étant donné que les exportations de service augmentent elles aussi de 50\% sur cette même période. Une hausse qui bénéficie à plusieurs industries clefs de l'UE comme l'automobile et l'électronique. Plusieurs sondages menés par l'EEAS soulignent d'ailleurs la bonne réception de cet accord par les entreprises européennes présentes en Corée. Trois quarts d'entre elles disent ainsi avoir étendue leur activité en ROC depuis l'accord et affichent un taux de satisfaction proche de 90\%.

Cependant, du côté coréen l'enthousiasme précédent l'accord fait vite place à un certain désenchantement face à des performances décevantes pour les exportations coréennes. De fait sur la même période les exportations de la Corée vers l'UE n'augmentent que de 32\%\textbf{{[}4{]}} alors que les importations en provenance d'Europe montent nettement plus rapidement, notamment dans les 5 ans après la signature de l'accord. Cette nouvelle tendance mène même à un inversement de la balance commerciale entre les deux parties en faveur de l'UE de 2011 à 2017 qui surprend les coréens auparavant en net excès. Les raisons avancées pour cette différence de performance sont nombreuses~: D'un point de vue conjoncturel, l'UE, plus encore que la Corée, souffre de l'impact de la crise financière de 2008 et le marché intérieur européen se révèle être nettement moins dynamique que prévu. Par ailleurs, les exportations de biens manufacturés coréens vers l'UE ont diminué. De fait, les industries coréennes qui exportent en Europe se réorganisent pour délocaliser la production en Europe là où auparavant celles-ci se limitaient à la distribution. Ainsi les IDE\textbf{{[}5{]}} coréen en direction des pays de l'Est de l'Europe comme la république tchèque augmentent et de grandes entreprises coréennes comme Kia motors implantent leurs premières usines en Europe pour assurer leur production destinée au marché intérieur européen. Enfin, au-delà de ce mouvement de délocalisation en Europe de l'Est, les entreprises coréennes délocalisent toujours plus en Asie du sud-Est, notamment au Vietnam. De fait, la Corée du Sud est le premier pays émetteur d'IDE en direction du Vietnam, avec plus de 19 milliards de dollars investis depuis 2008. Une tendance largement impulsée par les géants coréens de l'industrie, les chaebols, dont Samsung est l'un des plus importants. Aussi, une partie non négligeable des produits coréens sont exportés au Vietnam pour l'assemblage avant d'être réexporté en Europe ce qui fait que ces exports ne sont pas comptabilisés dans la balance EU-Corée.

Dès lors, sur le plan commercial cet accord demeure positif dans sa globalité et affiche des performances respectables, bien que décevantes comparé aux attendus coréens. Néanmoins la controverse autour de cet accord ne se concentre pas tant sur cet aspect commercial que sur d'autres aspects, d'avantage uniques à la politique commerciale européenne.

\partbreak

En premier lieu, l'insistance de l'UE à imposer à ses partenaires des normes fortes en matière sociale et environnementale est parfois mal reçu par ceux-ci. Plusieurs des principes forts de la charte de l'OIT qui n'avait pas été ratifié avant l'accord ne le sont toujours pas plusieurs années après. C'est notamment le cas de l'article portant sur la liberté d'association qui vise, entre autres, à faciliter la mise en place de syndicat, chose très difficile en Corée du Sud encore aujourd'hui. Aussi l'UE appelle le 5 juillet 2019 à ce qu'un panel d'experts soit constitué pour statuer sur le respect, ou non, par la Corée du Sud de ces principes. Ce panel d'expert penche en faveur de l'accusation de l'UE qui menace alors la Corée du Sud de possibles sanctions commerciales. Toujours dans le domaine des droits de l'homme, la législation autour de la peine de mort est aussi un important sujet de débat. Si la question est absente de l'accord, elle reste importante pour comprendre certaines des tensions qui subsistent entre l'UE et la Corée du Sud et impactent donc indirectement leurs relations commerciales. L'Europe a en effet appelé plusieurs fois la Corée à abolir la peine de mort, ingérence parfois mal reçue par une partie de l'opinion publique coréenne. On retrouve une tension similaire dans le cadre des négociations entre l'UE et l'Indonésie avec une perspective plus environnementale que sociale. En effet le ban européen sur l'huile de Palme fait craindre à l'Indonésie \textbf{{[}6{]}} une relation asymétrique à l'issue de cet accord et hésite donc à s'y engager comparé à d'autres accords moins exigeant environnementalement comme celui qu'elle a signé avec ses partenaires de l'ASEAN.

Par ailleurs, des barrières techniques et économiques subsistent encore entre les deux pays, en particulier autour de la «~Rule of Origin~». En effet si l'UE produit la majorité de ses biens d'exportations dans les limites du marché intérieur européen (notamment en délocalisant en Europe de l'Est pour les pays les plus riches), la Corée a souvent recours à la délocalisation dans des pays étrangers. De plus, certaines normes phytosanitaires font parfois débats comme celle portant sur la pasteurisation du fromage, procédé obligatoire en Corée mais pas toujours utilisé en Europe au profit d'autres techniques de production comme pour le parmesan qui est autorisé à l'export bien après le traité. Enfin le passage de déficit commercial rencontré par la Corée laisse un souvenir de l'accord très mitigé dans l'opinion publique coréenne. A titre d'exemple, en 2016, lors d'une conférence de presse organisée par l'UE pour célébrer le 5\textsuperscript{ème} anniversaire de l'ALE UE-Corée, aucun représentant du gouvernement coréen n'étaient présents.

En conclusion, à l'image de l'ALE entre l'UE et la Corée et des échecs de l'UE avec l'Indonésie, la volonté de l'UE de s'imposer comme une «~puissance normative~»\textbf{{[}7{]}} et commerciale se heurte parfois à des réalités culturelles et économiques antagonistes à la réalisation de ses objectifs. Cependant, c'est aussi ce qui fait la particularité de l'Union européenne~: son aspiration à diffuser un mode de développement qu'elle juge plus responsable, y compris par la voie commerciale qui semble pourtant ne pas toujours s'y prêter. S'il est encore difficile d'estimer précisément à l'heure actuelle dans quelle mesure cette approche freine la défense de ses intérêts en Asie, région que l'UE ne cesse de désigner comme une région stratégique sur le plan commercial et géopolitique, il est certain que c'est aussi une façon pour celle-ci de faire valoir ses valeurs dans un espace économiquement et idéologiquement dominé par les Etats-unis et la Chine.

\textbf{Notes~:}

\textbf{{[}1{]}} De fait, dans un rapport de l'EEAS du 21 février 2022, l'UE a présenté son plan pour «~ accélérer son engagement stratégique dans la région d'importance capitale qu'est l'indopacifique~». Cet engagement s'inscrit dans une dynamique de diversification des partenaires commerciaux de l'UE. En effet, l'UE a annoncé début juin 2022 vouloir redoubler d'effort pour signer de nouveaux ALE avec d'autres pays afin de réduire le problème de dépendance auquel elle fait face dans le cadre de la guerre entre l'Ukraine et la Russie.

\begin{itemize}
\item
  \textbf{Sources:}

  \begin{itemize}

  \item
    EU fact sheet on the EU strategy for cooperation the Indo-pacific, EEAS, 21 février 2022
  \item
    Philip Blenkinsop, EU seeks to speed up trade deals as Ukraine war forces rethink, 3 juin 2022, Reuters
  \end{itemize}
\end{itemize}

\textbf{{[}2{]}} Les barrières non tarifaires comprennent à titre d'exemple les politiques d'embargos, de licences ou encore de quotas.

\textbf{{[}3{]}} Les 4 axes de l'OIT que la Corée n'a pas encore ratifié sont les suivants~:

\begin{itemize}
\item
  Liberté d'association
\item
  Droit d'organisation et de négociation collective
\item
  Interdiction et abolition du travail forcé
\end{itemize}

\textbf{{[}4{]}} Ce chiffre correspond à l'évolution des exports de la Corée vers l'UE de 2010 à 2020. En réalité l'année après l'accord ceux-ci ont d'abord connu une baisse avant de remonter, notamment en 2017. Une baisse en partie compensé par l'effet d'anticipation de l'accord qui a causé une augmentation significative des exportations avant son entrée en vigueur

\textbf{{[}5{]}} Les IDE ou investissements directs à l'étranger correspondent «~aux investissements qu'une unité institutionnelle résidente d'une économie effectue dans le but d'acquérir un intérêt durable dans une unité institutionnelle résidente d'une autre économie et d'exercer une influence significative sur sa gestion dans le cadre d'une relation à long terme.~», définition donnée par l'INSEE, le 05/11/2019

\textbf{{[}6{]}} Selon des données publiées par Statista en avril 2022, l'Indonésie est le premier producteur mondiale d'huile de palme et assure jusqu'à 60\% de la production mondiale.

\begin{itemize}
\item
  Lien vers le rapport de Statista~: \url{https://fr.statista.com/infographie/27346/evolution-de-la-production-mondiale-huile-de-palme-et-principaux-pays-producteurs-en-2022/}
\end{itemize}

\textbf{{[}7{]}} Je reprends ici le terme utilisé par Zaki Laïdi, politologue français et directeur de recherche au CEVIPOF, dans son ouvrage \emph{EU foreign policy in a Globalized world}, qui l'oppose à des conceptions plus traditionnels de la puissance à l'image du «~hard power~» américain.

\textbf{Sources:}

\begin{itemize}
  \item William Yuen Yee, Can the EU and Indonesia Sign their elusive free trade agreement?, The Diplomat, 2 November 2022.
  \item Alexandra Hennessy and Poppy S. Winanti, Part IV, Chapter 14, EU-Indonesia Trade relations, A geo-economic turn in Trade policy? EU Trade Agreements in the Asia- Pacific Region, 2022, Palgrave Macmillan
  \item EU Strategy for Cooperation in the Indo-Pacific, European Union External Action, 21 February 2022
  \item Evaluation of the results of the EU-Korea Free Trade Agreement, European Union External Action, 2016
  \item Csilla Lakatos and Lars Nilsson, The EU-Korea Free Trade Agreement: anticipation, trade policy uncertainty and impact, note by the EU Chief Economist , June 2015.
  \item Dong-Hee Joe, A review of the First Decade of the Korea-EU FTA, World Economy Brief, vol. 12 No. 27, KIEP, 4 July 2022.
  \item Sunghoon Park, The new politics of trade negotiations: the case of the EU- Korea FTA, Journal of European Integration, 29 November 2017.
  \item Yoo-Duk Kang, Korea-Eu FTA: breaking new ground, Australia, the European Union and the New Trade agenda, 2017, ANU Press.
  \item Samuel Wejchert, EU-South Korea FTA: A ten-year perspective, European Institute for Asian studies, policy brief of 8 September 2021.
  \item European Commission website, section on EU trade relations with South Korea: facts, figures and latest developments.
  \item European Commission fact sheet on the EU-South Korea Free Trade Agreement: a quick reading guide
  \item European Commission fact sheet on the EU-South Korea Free Trade Agreement: 10 key benefits for the European Union
  \item Yoo-Duk Kang and Jun Yeup Kim, Trade and investment between Korea and EU after the Korea-EU FTA and its prospect, KIEP World economy update, vol.3 no.44, 27 September 2013.
  \item Hanns Günther Hilpert and Sunghoon Park, EU and South Korea: a model example for foreign trade policy relations, Part IV, Chapter 10, A geo-economic turn in Trade policy? The EU’s trade agreements in the Asia-Pacific region, 2022, Palgrave Macmillan.
  \item Zaki Laidi, EU foreign policy in a globalised world, normative power and social preferences, 2008.
  \item Trường Lăng, The Journey of Samsung in Vietnam, Viettonkin consulting, 26 July 2022.
  \item Jeon Jong-Hwi, South Korea to become a global test case in its trade dispute with the EU, Hankyoreh, 5 July 2019.
\end{itemize}

\end{document}
